% Generated draft: requires mathematical and source review.
\documentclass[11pt]{article}
\usepackage[margin=25mm]{geometry}
\usepackage{fontspec}
\setmainfont{Latin Modern Roman}
\usepackage{amsmath,amssymb,mathtools,unicode-math}
\setmathfont{Latin Modern Math}
\usepackage{enumitem,needspace,xurl,hyperref}
\setlength{\parindent}{0pt}
\setlength{\parskip}{5pt}
\setlength{\emergencystretch}{3em}
\newcommand{\N}{\mathbb{N}}
\newcommand{\Z}{\mathbb{Z}}
\newcommand{\Q}{\mathbb{Q}}
\newcommand{\R}{\mathbb{R}}
\newcommand{\C}{\mathbb{C}}
\newcommand{\F}{\mathbb{F}}
\newcommand{\A}{\mathbb{A}}
\newcommand{\PP}{\mathbb{P}}
\newcommand{\E}{\mathbb{E}}
\newcommand{\eps}{\varepsilon}
\newcommand{\dd}{\,\mathrm{d}}
\begin{document}
\section*{UnsolvedMath ID: 30005299}
\textbf{Source problem code:} catalogue entry
\subsection{Definitions and mathematical statement}
Let \(k\) denote the base field, which is not specified in the supplied record, and let \(V=k^4\), so that \(\PP(V)\cong\PP^3_k\). A quadric is the zero scheme of a homogeneous form in \(\operatorname{Sym}^2(V^*)\). There is an unresolved dimension convention: a linear system of projective dimension \(3\) corresponds to a vector subspace of dimension \(4\), whereas a net usually has projective dimension \(2\). The question and its supplied literature assessment use terminology that does not resolve this distinction.

Under the projective-dimension interpretation, write \(S=\PP(W)\), where \(W\subseteq\operatorname{Sym}^2(V^*)\) has vector-space dimension \(4\). Choose a basis \(q_0,q_1,q_2,q_3\) of \(W\) and coordinates \(x_0,x_1,x_2,x_3\) on \(V\). Define the derivative matrix and its degeneracy scheme by
\[
J_W(x)=\left(\frac{\partial q_i}{\partial x_j}(x)\right)_{0\leq i,j\leq3},\qquad
\mathcal W(W)=V\bigl(\det J_W(x)\bigr)\subseteq\PP(V).
\]
The zero scheme is independent of the chosen basis of \(W\). If the determinant is nonzero, it is a homogeneous quartic and defines a quartic hypersurface scheme. At a geometric point \([x]\), its vanishing means that some nonzero linear combination of the \(q_i\) has zero differential at \(x\). In characteristic different from \(2\), Euler's identity also makes that quadric vanish at \(x\), so this condition says that \([x]\) is a singular point of a member of the system. In characteristic \(2\), that equivalence requires further care.

This derivative-degeneracy construction is an explicit candidate interpretation of the Weddle surface in the question. The supplied excerpt does not itself define the Weddle construction, its scheme structure, or its treatment of base points. Consequently its agreement with the intended source definition requires mathematical review.

Under the usual unrestricted determinantal interpretation, a determinantal quartic is a quartic hypersurface admitting an equation
\[
F(x)=\det A(x),\qquad A(x)\in\operatorname{Mat}_{4\times4}(V^*),\qquad F\ne0.
\]
No symmetry of \(A\) is imposed in this interpretation. Let \(\mathcal D\) be the Zariski closure, in \(\PP(\operatorname{Sym}^4(V^*))\), of the classes of such nonzero determinants. The word ``general'' can then be interpreted as a property holding on a dense open subset of \(\mathcal D\). This choice of parameter locus is a proposed clarification, not a definition supplied by the record.

\textbf{Question, conditional on these conventions.} Characterize the quartic hypersurfaces \(X\subseteq\PP(V)\) for which there exists a subspace \(W\subseteq\operatorname{Sym}^2(V^*)\), with \(\dim_k W=4\) and \(\det J_W\ne0\), such that
\[
X=\mathcal W(W).
\]
In particular, does there exist a dense open subset \(U\subseteq\mathcal D\) such that every geometric point \([F]\in U\) admits such a representation
\[
[F]=[\det J_W]?
\]
The source does not specify the base field or characteristic, whether quartics must be integral or reduced, or whether equality is intended scheme-theoretically, set-theoretically, or up to projective equivalence. The displayed question uses hypersurface schemes in a fixed projective space solely to make this candidate formulation explicit; these choices remain unresolved source qualifications.

The source emphasizes that the interpolation matrix need not be symmetric but is constrained by a partially symmetric tensor. For the derivative construction above, its rows are gradients of quadrics, so its entries are not arbitrary independent linear forms. The source also mentions higher-degree and higher-dimensional extensions involving \(d-1\) tensor contractions. Its appended discussion connects generalized Weddle loci with a first-order analogue of zero-eigenvalue eigenvectors and with failures of the weak Lefschetz property in polynomial-ring quotients. It also discusses rectangular interpolation matrices, rank-drop dimensions, and a separate question about the \(d\)-Weddle locus of general finite point sets. These contextual statements do not supply additional hypotheses for the quartic question. The appended point counts and curve-degree formula are corrupted in the supplied transcription and cannot be stated precisely from this evidence.
\subsection{Short English statement}
The problem asks which quartic surfaces in \(\PP^3\) occur as Weddle surfaces of systems of quadrics, and specifically whether a general determinantal quartic does. Interpreting the system as having projective dimension \(3\) gives four independent quadrics and a quartic derivative determinant. The supplied assessment calls the problem open, but the field, dimension terminology, and exact Weddle convention remain unresolved. This is a statement-only draft requiring human mathematical review.
\subsection{Sources}
\begin{itemize}
\item Supplied local UnsolvedMath registry record, permanent ID \(30005299\), title ``Determinantal Quartics as Weddle Surfaces,'' source problem code \texttt{catalogue entry}; collection: Oberwolfach Reports Open Problems. No individual proposer is identified. The record supplies proposed year \(2022\).
\item Supplied primary-source citation: \emph{Mini-Workshop: Subvarieties in Projective Spaces and Their Projections} (2023), Oberwolfach Reports, pp.\ 3093--3119, DOI \texttt{10.4171/owr/2022/54}. \url{https://doi.org/10.4171/owr/2022/54}. The supplied extraction identifies the target as Question 2. The citation's year \(2023\), the proposed year \(2022\), and the DOI's \(2022\) component are retained without resolving their relationship.
\item Supplied literature source: Chiantini--Fagioli, \emph{Weddle loci of nets of quadrics}. \url{https://arxiv.org/abs/2510.16571}. Its contents and its relationship to the dimension convention above have not been independently verified here.
\item Supplied literature source: EMS Press, OWR 2022/54. \url{https://doi.org/10.4171/owr/2022/54}.
\item Literature assessment dated exactly \texttt{2026-08-22}: the supplied curation labels the problem open and states that the general determinantal-quartic characterization remains unresolved, while related work studies Weddle loci and tensor singularities without proving it. This is supplied dated curation, not an independently verified literature assessment.
\item The registry supplies status \texttt{open} and statement status \texttt{corrected\_verified}, with a note claiming comparison against the official report. These are supplied registry labels and verification claims; no source visit or independent verification is asserted by this draft.
\end{itemize}
\end{document}
